% mage-auto-heading: building-the-factory
% mage-auto-heading: the-feasibility-probe
% mage-auto-heading: the-two-rules-of-agentic-software-engineering
% mage-auto-heading: the-build-in-stages
% mage-auto-heading: the-build-measured-weekly
% mage-auto-heading: what-accumulated-and-why
% mage-auto-heading: inside-the-factory
% mage-auto-heading: what-the-factory-knows
% mage-auto-heading: the-joins-what-the-factory-can-derive-and-enforce
% mage-auto-heading: the-computation-graph-a-representation-that-took-over-execution
% mage-auto-heading: exploring-alternative-reasoning-engines
% mage-auto-heading: the-factory-in-operation
% mage-auto-heading: what-escalated-failures-became
% mage-auto-heading: the-supervisors-job
% mage-auto-heading: economics
% mage-auto-heading: liability
% Authoritative LaTeX manuscript; Markdown is a generated compatibility projection.

% <!-- part-title: Agentic Software Factories -->

% <!-- chapter-title: Inside DocAble's Software Factory -->

\chapter{Inside DocAble\textquotesingle s Software Factory}

% <!-- box-family: evidence -->

\begin{quote}
\mbox{}%
\subsection{This section illustrates}\label{heading:illustrates}

✓ The Modeling Principle · ✓ The Alignment Principle · ✓ Governance Conversion · ✓ Residual Human Judgment
\end{quote}

% <!-- point: one-engineer-directed-a-fleet | One engineer directed coding agents over the 31 complete weeks examined; the final snapshot contains 615,395 production lines, 2,333,471 apparatus lines, and 143 model modules. | terms: fleet-scale -->

DocAble was built by one engineer (me) directing a fleet of coding agents. The record examined here spans \magevalue{seip_weeks} complete weeks, from March 12 through October 14, 2026. At the final snapshot, the system contained \magevalue{seip_prod_loc} lines of production code and \magevalue{seip_apparatus_loc} lines of engineering apparatus, including tests, checks, tools, infrastructure, models, and indexed documentation. The history identified \magevalue{seip_models} model modules. These measurements describe the resulting software and its engineering environment; they do not measure how much useful work each line or model represents.

% <!-- point: the-engineers-role-was-engineering-not-implementation | New work did not originate in the agents or the factory's models. It came from engineering needs: goals, user feedback, production failures, standards obligations, advocacy guidance, and institutional use. The engineer decided which pressures became engineering work and made the consequential decisions about requirements, architecture, validation, and deployment. | terms: governed-engineering-environment -->

New work did not originate primarily in the coding agents or the factory\textquotesingle s existing models. It came from engineering needs: my own goals for the system, user feedback, failures observed in production, changes and obligations in accessibility standards, public guidance from accessibility advocacy groups,\citep{nfb-accessible-documents-2014} and the demands introduced as the system moved into institutional use. I decided which of these pressures became engineering work and made the consequential decisions about requirements, architecture, validation, and deployment.

% <!-- point: product-decisions-became-factory-obligations | Product requirements and architectural decisions became obligations for the environment in which coding agents changed DocAble. -->

I was engineering two coupled systems: DocAble and the factory that produced it. The product had to improve accessibility while preserving the document properties its users depended on. Those requirements shaped the architecture, and the factory had to preserve the resulting decisions as agents changed the implementation. When a library operation destroyed existing accessibility structure, for example, the product acquired a typed mutation interface, and the development environment acquired a check rejecting writes outside that interface. The repair changed both what the product could do and how future agents could change it.

% <!-- point: the-agents-performed-the-entire-implementation | The existing engineering environment shaped how the work could be realized: architecture, tests, models, instructions, validators, deployment machinery, and accumulated controls carried earlier decisions. Agents implemented inside it, and the engineer's attention went to direction, evidence, admission, and changing the production system. | terms: fleet-scale -->

The existing engineering environment then shaped how that work could be realized. Architecture, tests, system models, project instructions, validators, deployment machinery, and accumulated controls carried decisions made during earlier work. Agents performed most of the implementation inside that environment. I inspected almost none of the code they produced. My attention went instead to deciding what should change, resolving consequential design questions, examining evidence and failures, deciding what work could be admitted, and changing the production system when later work should inherit an earlier judgment.

% <!-- point: the-division-of-work-is-the-factory-examined | That division of work is the factory this section examines, in seven movements: the problem, the build with its longitudinal record, the mature factory's representations and machinery, its operation, the supervisor's changed job, the economics, and the liabilities that came with it. -->

That division of work is the factory examined in this section. The section moves in seven steps: the problem the factory was built for; the build, read from the repository\textquotesingle s own weekly record; the mature factory\textquotesingle s representations and machinery; that factory in operation; the supervisor\textquotesingle s changed job; the economics; and the liabilities that came with it. The chronology matters because none of it arrived fully formed. As realization accelerated, failures exposed missing representations, controls, and supervisory surfaces; the production system changed in response.

% <!-- point: requirements-ordinary-capacity-unusual | The requirements were ordinary production-software requirements; the unusual condition was the implementation capacity available to one engineer. | terms: governance-target-product -->

The requirements were not unusual for production software. \magevalue{chapter:problem} introduced them by kind; this section states them in full, where they begin to do their work. What was unusual was the implementation capacity available to one engineer.

\section{The Problem}\label{heading:nearly-free}

% <!-- section-terms: doj-2024-rule, remediation-cost -->

% <!-- point: the-rule-and-the-extension | The 2024 DOJ rule set WCAG 2.1 AA and a 2026 deadline; four days before the deadline, the Department extended it by a year. | terms: doj-2024-rule, ada-title-ii -->

In 2024, the U.S. Department of Justice required state and local governments, including public universities and schools, to make covered web content and mobile applications conform to \magevalue{doj_technical_standard}.\citep{doj-web-accessibility-2024} The original compliance date for larger public entities was April 24, 2026. Four days before that deadline, the Department extended it by one year.\citep{doj2026extension}

% <!-- point: the-explanation-was-a-production-problem | The Department's explanation described a production problem: technology, including generative AI, had not delivered remediation at scale. | terms: doj-2024-rule -->

The Department\textquotesingle s explanation described a production problem. Public entities held large amounts of inaccessible material and limited staff to remediate it. Correspondence cited in the rule argued that current technology, "including generative AI," could not reliably automate the remediation of STEM materials at scale. The Department itself concluded that the advancement and availability of technology "did not meet the Department\textquotesingle s expectations."

% <!-- point: institutional-remediation-made-the-workload-concrete | The undergraduate effort motivates requirements for scale and reviewable quality; its measured outcomes appear in Economics. -->

The undergraduate remediation effort introduced in \magevalue{sec:mage-by-example} made the workload concrete: even trained staff needed substantial time to repair ordinary course materials. Automation would have to reduce that recurring effort while producing results instructors could review and trust. The comparison later in this section examines both the outputs and the remaining limits of the evidence.

% <!-- point: box-checking-sharpened-the-target | The Department recorded a quality concern — that the workload would produce procedural box-checking rather than sustainable accessibility work — and that concern sharpened the target: passing a checker is neither making a document accessible nor producing work an institution could defend in a lawsuit. | terms: doj-2024-rule, wcag-conformance -->

The Department also recorded a quality concern: institutions facing this workload might respond with rapid procedural box-checking rather than sustainable accessibility work. Passing an accessibility checker is not the same as making a document accessible --- or producing work an institution could defend in a lawsuit.\magefootnote{accessibility-litigation}{Digital-accessibility litigation is tracked continuously rather than reported as a closed figure. UsableNet\textquotesingle s ADA Accessibility Lawsuit Tracker enumerates federal and key state-court digital-accessibility cases on a monthly cadence, and recorded 4,928 web-accessibility lawsuits filed during 2025.\citep{usablenet-lawsuit-tracker} Laura Carlson maintains the higher-education-specific tally of lawsuits, complaints, and settlements.\citep{CarlsonHigherEdAccessibilityLawsuits} Both are living sources; the count above is the reading at the access date recorded in the bibliography, not a settled total.}

% <!-- point: the-eight-requirements-stated-where-they-work | The eight requirements are stated in full here, where they do their work: {{sec:mage-by-example}} grouped them into kinds of obligation to establish the problem; stated in full, they explain why demonstrating semantic reasoning was only the beginning. | terms: governance-target-product -->

\magevalue{sec:mage-by-example} grouped these obligations by kind to establish the nature of the problem. \autoref{docable-requirements} states the eight requirements: they explain why demonstrating semantic reasoning was only the beginning.

{\footnotesize
\begin{longtable}[]{@{}
  >{\raggedright\arraybackslash}p{(\linewidth - 2\tabcolsep) * \real{0.5000}}
  >{\raggedright\arraybackslash}p{(\linewidth - 2\tabcolsep) * \real{0.5000}}@{}}
\caption[The DocAble requirements]{What DocAble had to accomplish. The eight requirements, R1--R8 --- the specification the factory had to satisfy, which \magevalue{sec:mage-by-example} introduced by kind.}\label{docable-requirements}\tabularnewline
\toprule\noalign{}
\begin{minipage}[b]{\linewidth}\raggedright
Requirement
\end{minipage} & \begin{minipage}[b]{\linewidth}\raggedright
DocAble must\ldots{}
\end{minipage} \\
\midrule\noalign{}
\endfirsthead
\toprule\noalign{}
\begin{minipage}[b]{\linewidth}\raggedright
Requirement
\end{minipage} & \begin{minipage}[b]{\linewidth}\raggedright
DocAble must\ldots{}
\end{minipage} \\
\midrule\noalign{}
\endhead
\bottomrule\noalign{}
\endlastfoot
R1 · Multi-format remediation & Remediate educational documents in the formats people actually use, including Office (PowerPoint, Word, Excel) and PDF. \\
R2 · Semantic remediation & Perform semantic repairs that require understanding document content, not only syntactic, deterministic transformations. \\
R3 · Preserve the document & Retain content, appearance, and behavior except where remediation requires a change. \\
R4 · Standards-grounded accessibility & Evaluate and repair documents against their governing obligations (WCAG 2.1 AA), not proxy checker scores. \\
R5 · Evidence & Record what accessibility requirement motivated each change and what evidence supports the resulting document state. \\
R6 · Reviewability & Make outputs auditably traceable: consequential automated changes are inspectable and, where appropriate, reversible. \\
R7 · Production operation & Process real documents at practical sizes, costs, and latency. \\
R8 · Institutional deployment & Operate as a deployed service that scales to thousands of faculty at one institution, and preferably nationwide. \\
\end{longtable}
}

% <!-- point: requirements-explain-the-size | Semantic reasoning was only one requirement of eight. The capability alone did not establish preservation, reviewable changes, conformance, or dependable service operation. Those obligations generated the engineering system around the model. | terms: governance-target-product -->

These requirements explain why the resulting system grew large. Semantic reasoning was only one requirement of eight. A frontier model could supply a capability that had previously been missing. That capability alone did not establish preservation, reviewable changes, conformance, or dependable service operation; each remained an engineering obligation. Those obligations generated the engineering system around the model.

% <!-- point: the-first-question-was-whether-the-capability-was-real | Seven of the eight requirements were engineering. The first question was whether frontier models could supply the missing capability. | terms: feasibility-probe -->

The first question was whether frontier models could supply the missing capability.

\section{Building the Factory}\label{heading:building-the-factory}

% <!-- section-terms: feasibility-probe, stage-sequence, support-ratio -->

\subsection{The Feasibility Probe}\label{heading:the-feasibility-probe}

% <!-- point: docable-began-with-a-feasibility-question | DocAble began with a feasibility question, because remediation requires semantic reasoning: figures, tables, reading order, hierarchy, emphasis, and visual arrangement all carry meaning a machine must recover to present through another channel. | terms: feasibility-probe, semantic-artifact -->

% <!-- index-def: semantic-artifact -->

DocAble began with a feasibility question. Document remediation requires semantic reasoning. A document is not merely a container of text: figures, tables, reading order, hierarchy, emphasis, and visual arrangement all carry meaning. Making that document accessible requires recovering enough of that meaning to present it through another channel --- to a screen reader, to a keyboard, or to a student who cannot see the figure a lecturer drew.

% <!-- point: conventional-automation-could-not-supply-that-reasoning | Automated accessibility had historically been unable to supply that reasoning — too many formats, and too many ways of expressing the same idea, for conventional automation to enumerate — so modern frontier models were obvious candidates. | terms: semantic-gap, vision-language-model -->

Automated accessibility had historically been unable to supply much of this reasoning. Documents come in too many formats, and humans express the same ideas in too many ways, for conventional automation to enumerate the possibilities. Earlier reasoning systems could not reliably bridge that semantic gap. Modern frontier models were obvious candidates.

% <!-- point: the-probe-tested-the-uncertain-part-first | The probe gave a multimodal model pictures of equations, technical slides, and handwritten notes; it interpreted them and produced useful transcriptions and alternative text. | terms: feasibility-probe, vision-language-model -->

I gave a multimodal model pictures of mathematical equations, technical slides, and handwritten notes. It had no trouble interpreting them and producing useful transcriptions and alternative text.

% <!-- point: the-rest-was-engineering | The probe established that generative AI could supply the semantic reasoning remediation needed; building a reliable production system around it remained an engineering problem. | terms: feasibility-probe, semantic-gap -->

The feasibility probe established that generative AI could supply the semantic reasoning needed for remediation. Building a reliable production system around that capability remained an engineering problem.

% <!-- point: engineering-set-the-price-of-the-reasoning | Different structure-inference experiments illustrate how an engineered representation changes the reasoning task; unpaired fixtures do not establish an equivalent-quality speedup. | terms: feasibility-probe -->

Later experiments examined what an engineered representation could contribute. In one setting, frontier agents inferred PDF structure from scratch through multiple tool interactions. In another, a frontier model inferred structure in one pass over a prepared scaffold on a separate clean-room prose fixture. The scaffolded pass achieved high agreement with that fixture's reference structure at much lower inference cost; cheaper models could also use the scaffold. The experiments involved different documents and interaction patterns, so they do not establish a matched cost reduction at equivalent quality. They do illustrate an engineering choice: what the model must reconstruct and what useful structure it can inherit. The whole-document agent comparison in the Economics section provides a more direct, though still small, comparison of remediation workflows.\citep{davis2026docable}

% <!-- point: cost-was-only-one-reason-not-to-hand-the-document-to-the-model | Cost was one reason among several: DocAble separated semantic judgment from consequential mutation, so models reasoned over representations, typed operations changed the artifact, independent validators checked mechanically establishable properties, and provenance was embedded with mutations so the change record could be reconstructed from the document. -->

Cost was only one reason not to make the model the remediation system. DocAble instead separated semantic judgment from consequential mutation. Models reasoned over representations; typed operations changed the artifact; independent validators checked properties that could be established mechanically; and provenance was embedded with consequential mutations, so that the eventual change record could be reconstructed from the document rather than trusted as a model-generated account.

\subsection{The Two Rules of Agentic Software Engineering}\label{heading:the-two-rules-of-agentic-software-engineering}

% <!-- point: not-reading-the-code-was-deliberate | Not reading the code was deliberate: inspecting agent-produced implementation would have bounded implementation throughput by the engineer's own reading capacity, so the first rule made reading the last resort and an engineering environment the surface through which work was directed and judged. | terms: fleet-scale -->

I began with two rules.

\textbf{Reading the code is the last resort.} Not reading the code was deliberate. If building DocAble required me to inspect the implementation produced by coding agents, then implementation throughput would remain bounded by my capacity to read code. The point was not to become faster at code review. It was to build an engineering environment in which I could direct work, understand the system, evaluate evidence, and decide whether changes were acceptable without routinely reconstructing those answers from implementation.

% <!-- point: reading-the-code-was-the-escape-hatch | The rule did not forbid reading source: needing to read it pointed to a missing supervisory surface — a model, test, lint, instrument, or architectural boundary — so reading was the escape hatch, not the operating model. | terms: fleet-scale -->

This did not mean that I refused to read source code. It meant that needing to read it pointed to a missing supervisory surface: a model, test, lint, instrument, architectural boundary, or other representation of the property I actually needed to reason about. Reading the code was the escape hatch, not the operating model.

% <!-- point: good-fast-and-cheap-became-a-rule-for-the-factory | The second rule required the resulting system to be good, fast, and cheap — partly a DocAble requirement, partly a test of whether agents could help build a production system good enough to trust, fast enough to matter, and cheap enough to deploy at scale. -->

\textbf{The resulting system should be good, fast, and cheap.} This was partly a requirement for DocAble and partly a test of agentic software engineering itself. I wanted to know whether agents could help build a production system that was not merely fast to develop, but good enough to trust, fast enough to matter, and cheap enough to deploy at scale. I therefore made all three requirements of the factory.

% <!-- point: failing-any-one-of-the-three-fails-the-test | A fast factory producing poor software would fail the test, and so would a high-quality system that stayed prohibitively expensive to develop or operate; agentic implementation was valuable only if the engineering system around it achieved all three. -->

A fast factory producing poor software would fail the test. So would a high-quality system whose development or operation remained prohibitively expensive. Agentic implementation was valuable only if the engineering system around it could achieve all three.

\subsection{The Build, in Stages}\label{heading:the-build-in-stages}

% <!-- point: the-build-moved-under-concrete-pressures | The build moved chronologically under concrete pressures: PowerPoint first, then more formats, deployment, real users, inadequate checkers, and hardening. Most new work entered through pressures outside the fabricator, which the engineer translated into engineering work that the accumulated production system then shaped. | terms: stage-sequence -->

I opened an empty repository and began with PowerPoint. The first program walked the contents of a slide deck and sent figures and equations to a model for description. Colleagues asked whether it could handle PDFs, so I generalized the system to Word and Excel and began adding PDF support (R1); documents from other people immediately exposed assumptions my own files had not. Putting the system online --- \href{\%7B\%7Bproduct_url\%7D\%7D}{DocAble} runs in beta at \href{\%7B\%7Bproduct_url\%7D\%7D}{scholaccess.com} --- introduced authentication, deployment, institutional use, and feedback from documents I had never seen (R8). Users then found outputs that passed existing accessibility checkers while losing semantic equivalence, so I replaced those checkers as the primary validation mechanism with validation grounded directly in the accessibility standards (R4). When the Department of Justice extended the compliance deadline, the immediate deadline pressure receded, and I shifted the project toward production hardening (R7). Most new work entered the factory through pressures outside the fabricator: new users and documents, observed failures, deployment demands, standards obligations, and my own engineering objectives. I translated those pressures into engineering work; the accumulated production system shaped how agents could realize it. \autoref{build-stages-timeline} gives the sequence. In a few months, a feasibility experiment had become a production software system.

\begin{figure}[htbp]
\centering
\includegraphics[width=\linewidth]{book/latex/assets/build-stages-timeline-a24e34269b.pdf}
\caption{\emph{The Seven Build Stages.} Retrospective sequence from the March feasibility probe to the July--September period of deployment, iteration, and optimization, giving the dates, the pressure, and the major engineering consequence at each stage. The serverless re-platforming is one episode inside that final period rather than a stage of its own.}
\label{build-stages-timeline}
\end{figure}

% <!-- point: what-the-engineer-did-while-agents-implemented | The engineer's role, visible in the stage narrative's own decisions: deciding what should change and on what evidence — retiring checker scores for standards-grounded validation, reading the deadline extension as a shift in engineering risk — while agents implemented, and changing the environment when a judgment should not recur. -->

The stages also show how my work differed from the agents\textquotesingle{} work. Take the checker replacement apart: the decision --- retire checker scores as the primary validation and ground validation in the accessibility standards instead --- was mine; the implementation was the agents\textquotesingle; the evidence that forced the choice was the users\textquotesingle{} documents, which had passed the checkers and still lost meaning. The response to the deadline extension separates the same way: reading the extension as a shift in engineering risk, away from shipping against a date and toward production hardening, was my decision; the hardening work it implied was the agents\textquotesingle. That division held throughout: I directed consequential changes and established what evidence would justify accepting them; agents implemented and investigated; and when a failure exposed a weakness in the system or its engineering environment, I changed the environment so the same judgment would not have to be made again.

% <!-- point: format-breadth-had-two-justifications | Supporting another document format had two distinct justifications: product value — users need the format — and information value: implementing it tests whether the architecture is general. | terms: governance-target-product -->

One class of decisions concerned which formats to support. Multi-format remediation was already a product requirement: users needed Office and PDF. But I deliberately supported additional formats, including LibreOffice documents and editable sources such as LaTeX and Typst, partly for what implementing them could teach us about the architecture.

If code is cheap and abstractions are expensive, how do we discover the right abstractions? One way is to use implementation to test them. An architecture developed around PowerPoint may appear general while quietly depending on its conventions. Supporting spreadsheets, other document formats, and editable source representations tests whether the interfaces capture the underlying remediation problem or merely the first formats implemented.

% <!-- point: breadth-tests-depth | The technique: propose an abstraction, implement additional members of the format family, and examine where special cases appear — essential format differences separated from accidental assumptions, with the engineer deciding whether to revise. -->

The technique is to propose an abstraction, implement additional members of the family, and examine where special cases appear. Some differences are essential to the formats; others reveal accidental assumptions in the architecture. The engineer must distinguish between them and decide whether to revise the abstraction. The additional implementations therefore serve two purposes: delivering useful capabilities and providing information about the design. When implementation is inexpensive, use implementation breadth to test abstraction depth.

\subsection{The Build, Measured Weekly}\label{heading:the-build-measured-weekly}

% <!-- point: four-traces-read-in-sequence | The SEIP reconstruction covers 31 complete weeks; work, recorded completions, controls, and source composition describe different aspects of the factory. -->

The repository record lets us examine the build beyond the retrospective stage narrative. The SEIP reconstruction covers \magevalue{seip_weeks} complete seven-day weeks from March 12 through October 14, 2026.\citep{davis2026docable} Four traces describe different aspects of the factory: landed work, documented completions and dispatches, accumulated models and verifiers, and the source supporting production. None alone measures productivity.

The work-commit series in \autoref{realization-over-time} reconstructs \magevalue{seip_commits} landed commits. It expands merge-train squashes into their source commits, deduplicates them, and excludes empty commits and integration bookkeeping. This matters because a change in landing mechanism would otherwise look like a change in engineering activity. Some source objects remain unavailable, so the reconstruction is incomplete; a commit still says little about the value or correctness of the work it records.

\begin{figure}[htbp]
\centering
\includegraphics[width=\linewidth]{book/latex/assets/docable-seip-work-261028.pdf}
\caption{\emph{Realization Over Time.} Reconstructed landed work commits in each complete seven-day project week, beginning March 12, 2026. Squash source commits are expanded and deduplicated; bookkeeping and empty commits are excluded. These counts measure recorded work, not product value or productivity.}
\label{realization-over-time}
\end{figure}

An Epic represented formalized engineering intent: a substantial desired change with a rationale, a design, relevant models and invariants, and phased acceptance criteria. Design came before implementation; nontrivial designs received review by another model, and a Definition-of-Done review assessed completion. One Epic could require several dispatches and many commits. The revised mining pass counts dated affirmative completion claims rather than moves into a closed directory. It recovers \magevalue{seip_completions} completions and \magevalue{seip_dispatches} agent dispatches in the complete window (\autoref{work-over-time}). These are records of activity and claimed completion, not independent assessments of success. Early missing records and uneven dispatch retention prevent treating an empty week as no work or comparing weekly dispatch rates as if coverage were uniform.

\begin{figure}[htbp]
\centering
\includegraphics[width=\linewidth]{book/latex/assets/docable-seip-activity-261028.pdf}
\caption{\emph{Recorded Work Over Time.} Documented Epic completions and recovered agent dispatches across the complete project weeks. Missing early records do not establish an absence of work. A completion is a recorded claim, not an independently verified outcome.}
\label{work-over-time}
\end{figure}

Week 11 marks a deliberate change in how I governed that work. My contemporaneous notes describe moving repeated review obligations from probabilistic review into deterministic checks, while commit activity and the formalization of Epic records show the environment changing around that decision. The convergence supports the interpretation of a change in engineering practice. The first surviving Epic or dispatch record, however, dates a recording convention rather than proving that the underlying practice began that week.

\subsection{The Controls Accumulate}\label{heading:control-growth}

% <!-- point: the-control-count-is-evidence-of-investment | Models, multi-model consumers, and registered verifiers accumulated; their counts describe structure, not its quality. | terms: governance-conversion -->

The third trace shows the structure that accumulated around implementation (\autoref{control-growth-curve}). At the last complete week, the reconstruction identifies \magevalue{seip_models} declared model modules, \magevalue{seip_consumers} consumers of at least two models, and \magevalue{seip_static} static and \magevalue{seip_dynamic} dynamic verifier files. Model existence is traced backward by path to avoid mistaking a later declaration convention for the moment the model was created. The verifier census counts existing files referenced by the registries; it does not count all tests or establish that each verifier ran. The engineering significance lies in what the structures represent and check, which the following examples examine.

\begin{figure}[htbp]
\centering
\includegraphics[width=\linewidth]{book/latex/assets/docable-seip-analyses-261028.pdf}
\caption{\emph{Models and Verifiers Accumulate.} Declared model modules, consumers of multiple models, and registered static and dynamic verifier files at weekly snapshots. Counts describe represented structure and available checks, not their correctness or execution.}
\label{control-growth-curve}
\end{figure}

\subsection{The Support Ratio}\label{heading:support-ratio-curve}

% <!-- index-def: support-ratio -->
% <!-- point: the-environment-outgrew-the-product | At ten consistently categorized revisions, supporting apparatus grew to several times production source; the ratio measures investment rather than its return. | terms: support-ratio, governance-target-product -->

The fourth trace measures how much supporting source accumulated relative to production source (\autoref{support-ratio-curve}). Production includes remediation, web processing, frontend, and shared services. The supporting apparatus includes tests, lints, tools, agent infrastructure, infrastructure-as-code, model metadata, and indexed documentation. Historical snapshots and generated copies are excluded. These are code-line counts at ten sampled revisions, not weekly newline counts over the entire repository.

Call apparatus divided by production the \textbf{support ratio}. It rises above parity early, reaches about 4.5 at week 17, and ends at \magevalue{seip_support_ratio}: \magevalue{seip_apparatus_loc} apparatus lines supporting \magevalue{seip_prod_loc} production lines. The composition matters as much as the ratio: tests, models, and infrastructure carry different obligations. Neither more source nor a higher ratio is a target. Supporting source earns the name engineering capital only while future work inherits useful capacity from it; the figure measures what accumulated, not its return.

\begin{figure}[htbp]
\centering
\includegraphics[width=\linewidth]{book/latex/assets/docable-seip-capital-261028.pdf}
\caption{\emph{The Product and Its Supporting Apparatus.} Source composition and the apparatus-to-production ratio at ten sampled revisions through October 14. Lines between samples interpolate; they are not additional measurements. The apparatus includes tests, lints, tools, agent infrastructure, infrastructure-as-code, model metadata, and indexed documentation.}
\label{support-ratio-curve}
\end{figure}

Together, the traces show sustained realization accompanied by an expanding engineering environment. They do not establish that a particular amount of modeling caused a particular increase in output. The rest of the case examines the mechanisms behind that accumulation: what decisions the structures preserved, which failures they addressed, and what work remained with the engineer.

\subsection{What Accumulated, and Why}\label{heading:what-accumulated-and-why}

% <!-- point: ordinary-and-agent-specific-changes | What accumulated was a factory around the implementation: some of it ordinary software engineering, other parts existing specifically because agents did the implementation. -->

What had accumulated? Not only implementation, but a factory around it. Some of it was ordinary software engineering. Tests accumulated. Static analyses checked properties that mattered. Deployment gates became stronger. Architectural boundaries became explicit. Other parts existed specifically because agents were doing the implementation: task briefs, agent-facing project instructions, role-specific dispatch, mechanisms for selecting the context relevant to a task, and controls over concurrent work.

% <!-- point: model-selection-was-an-engineering-decision | Model selection was itself an engineering decision: frontier models were reserved for judgments where their additional capability earned their additional cost — the cheapest-adequate-judge rule applied to the fabricator rather than to the validator. -->

Model selection was itself an engineering decision: frontier models were reserved for judgments where their additional capability earned their additional cost --- the cheapest-adequate-judge rule of \magevalue{sec:engineering-the-environment} applied to the fabricator rather than to the validator.

% <!-- point: parallel-implementation-cheap-admission-serialized | The resulting arrangement made parallel implementation cheap while admission to the production system remained serialized through evidence: realization capacity fanned out without granting each fabricator authority to admit its work. -->

The resulting arrangement made parallel implementation cheap while admission to the production system remained serialized through evidence. Realization capacity could fan out without granting each fabricator independent authority to admit its work.

% <!-- point: the-two-sides-interacted | The distinction was not clean: the engineering environment and the product developed together. -->

The distinction was not clean, because the two sides interacted. A model of the software architecture could guide an agent before it worked and support a deterministic check afterward. A failure discovered by an agent could expose an architectural weakness whose repair also made future agent work easier to control. The engineering environment and the product developed together.

% <!-- point: the-environment-accumulated-in-sequence | The environment was absent at the start; April, May, and July each added durable structure in response to failures. | terms: governance-conversion -->

And none of it was present when I started. In April, a production failure led me to introduce an explicit state model for a processing lifecycle. In May, duplicated descriptions of the system were consolidated into a shared modeling substrate. By July, properties represented in those models were determining what verification the repository required. Later failures led to additional models and controls for service interactions, deployment, entitlements, and resource use.

% <!-- point: the-sequence-matters | The work began with a software problem and capable agents, not with MAGE or a plan for a factory; the structures appeared at the limits of directing agents without them. | terms: governance-conversion -->

The sequence matters because I did not begin with MAGE, or with a plan for an agentic software factory. I began with a software problem and capable implementation agents. The engineering structures described in this chapter appeared as I encountered the limits of directing those agents without them.

% <!-- point: the-exhibits-measure-what-the-requirements-cost | The exhibits record implementation and supporting structure accumulated while addressing the requirements, not the cost or causal contribution of each requirement. | terms: governance-target-product -->

The exhibits record the implementation and supporting structure that accumulated while I addressed these requirements. They do not isolate the cost of any one requirement. The following examples show how particular obligations shaped that structure.

% <!-- box-family: evidence -->

\begin{quote}
\mbox{}%
\subsection{Field note --- The counterfactual}\label{heading:fieldnote-counterfactual}

One counterfactual rests on direct engineering judgment rather than repository data. \textbf{Within the actual constraints of this project---one primary engineer, an academic budget, and the development interval reported here---DocAble was not feasible for me without generative AI.} That is not a claim that a conventional engineering organization could not have built the system. The relevant counterfactual is the project that actually existed: under those constraints, I would not have attempted a production system at this scope without commodity implementation capacity.

A second counterfactual concerns sustained delegation: \textbf{GenAI made DocAble feasible; the machinery that became MAGE made delegation at this rate feasible.} Once autonomous change volume exceeded anything I could inspect conventionally, raw model capability was no longer enough. Continuing at that rate required increasingly explicit representations of the system, independent evidence, constraints, validators, gates, and operating machinery through which the work could be understood and governed.

This is a within-case engineering judgment, not a controlled causal estimate.
\end{quote}

% <!-- point: the-facts-reinterpreted-around-supervisory-capacity | DocAble did not solve abundant implementation by reading faster; it changed what the engineer had to supervise. | terms: governed-engineering-environment -->

DocAble did not solve abundant implementation by finding a way for its engineer to read implementation faster. It changed \textbf{what the engineer had to supervise}. As implementation accelerated, the engineering environment grew to include representations of the system, controls over implementation, and mechanisms for converting recurring failures into durable constraints. The next sections examine the factory\textquotesingle s representations and operation, followed by the changes in the engineer\textquotesingle s supervisory role.

\section{Inside the Factory}\label{heading:inside-the-factory}

% <!-- section-terms: modeling-principle, alignment-principle, governed-engineering-environment -->

% <!-- point: source-code-stopped-being-an-adequate-interface | Source code alone stopped being an adequate engineering interface as DocAble grew: repository search could locate a symbol but not say which services participated in a workflow, which component owned a piece of state, or what had to stay bounded as a path changed. | terms: modeling-principle -->

Source code alone stopped being an adequate engineering interface as DocAble grew. Repository search could locate a symbol; it could not say which services participated in a workflow, which component owned a piece of state, or what had to stay bounded as a path changed. This subsection walks the interior of the factory the weekly series just measured: the representations the mature factory maintains, the joins that connect them, and the one representation that took over execution.

\subsection{What the Factory Knows}\label{heading:what-the-factory-knows}

% <!-- point: the-census-is-live-not-frozen | The 31-week reconstruction identifies 143 model modules; models remain maintained artifacts rather than a fixed inventory. | terms: modeling-principle -->

The 31-week reconstruction identifies \magevalue{seip_models} model modules at the final snapshot. These representations change with the system and answer different engineering questions. Registries describe ownership and available operations; budgets make resource bounds explicit; state machines expose legal transitions; decision tables preserve policy distinctions. Some models describe the product, while others describe the agent fleet and its engineering environment. The factory represents both the software it produces and the process that produces it.

% <!-- point: specimens-show-why-representations-afford-supervision | Six specimens show the models' purpose: each representation answers an engineering question for a consumer, and its supervisory consequence is what the engineer no longer reconstructs from implementation. | terms: modeling-principle -->

The point is what those models afford. Chapters \magevalue{chapter:modeling|num}--\magevalue{chapter:one-problem-many-models|num} argued that a model is a purposeful reduction matched to an engineering question. What the earlier chapters could not yet show is the column a factory adds: the \textbf{supervisory consequence} --- what the engineer stops having to reconstruct from implementation because the representation carries it. \autoref{factory-census-specimens} shows six representative models, chosen for the variety of question, form, and consumer.

{\footnotesize
\begin{longtable}[]{@{}
  >{\raggedright\arraybackslash}p{(\linewidth - 6\tabcolsep) * \real{0.2500}}
  >{\raggedright\arraybackslash}p{(\linewidth - 6\tabcolsep) * \real{0.2500}}
  >{\raggedright\arraybackslash}p{(\linewidth - 6\tabcolsep) * \real{0.2500}}
  >{\raggedright\arraybackslash}p{(\linewidth - 6\tabcolsep) * \real{0.2500}}@{}}
\caption[Census specimens and their supervisory consequences]{Six representative engineering models. Each row: the engineering question, the representation that answers it, what consumes it, and what the supervisor no longer reconstructs from implementation.}\label{factory-census-specimens}\tabularnewline
\toprule\noalign{}
\begin{minipage}[b]{\linewidth}\raggedright
Engineering question
\end{minipage} & \begin{minipage}[b]{\linewidth}\raggedright
Representation
\end{minipage} & \begin{minipage}[b]{\linewidth}\raggedright
Consumer
\end{minipage} & \begin{minipage}[b]{\linewidth}\raggedright
Supervisory consequence
\end{minipage} \\
\midrule\noalign{}
\endfirsthead
\toprule\noalign{}
\begin{minipage}[b]{\linewidth}\raggedright
Engineering question
\end{minipage} & \begin{minipage}[b]{\linewidth}\raggedright
Representation
\end{minipage} & \begin{minipage}[b]{\linewidth}\raggedright
Consumer
\end{minipage} & \begin{minipage}[b]{\linewidth}\raggedright
Supervisory consequence
\end{minipage} \\
\midrule\noalign{}
\endhead
\bottomrule\noalign{}
\endlastfoot
Can a job always terminate legally? & Composed cross-service state-machine model with per-invariant facts & Exhaustive-state checkers; the workers\textquotesingle{} own transition tables & Legal progress and termination are reasoned over states and transitions, not a multi-thousand-line worker \\
How much memory can a job need? & Peak-memory sizing model that upper-bounds a job and derives the deployment memory tier & Deployment configuration; regression tests on the bound & A resource bound is inspected directly rather than reconstructed from allocation behavior \\
Which services exist, and what may call what? & Deployment topology plus a typed inter-service call-edge model & Deploy machinery; agents localizing changes & Ownership and boundaries are read from the map, not rediscovered per change \\
May this user submit this document? & Entitlement decision table --- the single policy authority for the admission decision & The front door at upload time & The admission rule is one table; changing policy is editing a model, not hunting scattered conditionals \\
What may touch a document format\textquotesingle s internals? & Registry of external seams plus typed format models owning all mutation & Build-time ban checks on raw library access & A format-corruption class is excluded by construction rather than hunted after release \\
Which step changed this document element? & Per-mutation provenance stamps projected into a derived changelog & The customer-facing changelog; forensic queries & Diagnosis is a lookup over evidence, not a cross-file investigation \\
\end{longtable}
}

% <!-- point: the-supervisor-does-not-supervise-2mloc-through-2mloc | The recurring pattern: each representation replaces implementation-reading with a purpose-built supervisory surface — supervision runs through representations sized to the engineering question. | terms: modeling-principle -->

Walk any row and the same pattern appears. The lifecycle model lets its supervisor reason about legal progress and termination rather than reading the worker that implements them; the memory budget lets the supervisor reason about a bound rather than reconstructing allocation behavior; the entitlement table preserves an admission-policy distinction as one authoritative representation rather than logic diffused through handlers. The implementation and its apparatus contain millions of lines. \textbf{Supervision does not require reading all of those lines} --- supervision runs through representations sized to the questions that matter.

\subsection{The Joins: What the Factory Can Derive and Enforce}\label{heading:the-joins-what-the-factory-can-derive-and-enforce}

% <!-- point: a-repository-of-documentation-is-not-a-supervisory-environment | Models alone would be documentation; the factory joins them and gives them consequences, which is the difference between describing knowledge and acting on it. | terms: alignment-principle -->

A repository can contain excellent documentation and still supervise nothing. What separates DocAble\textquotesingle s models from documentation is that the environment \textbf{joins} them and \textbf{acts} on them: representations connect to other representations, and selected relations carry mechanical consequences. The architecture resembles proof checking in one limited respect: a powerful producer is separated from the smaller mechanism that decides whether its evidence is sufficient for admission.\citep{ringer2019qed} Three joins show the difference in increasing order of reach.

% <!-- point: the-invariant-chain-derives-checker-obligations | The invariant chain: property facts on each invariant derive a verification tier, the tier mandates a checker class, and a check fails when the checker is missing — the model tells the factory what evidence it owes. | terms: alignment-principle -->

\textbf{Property facts → verification tier → required checker → blocking consequence.} Each invariant in the cross-service state-machine model carries declared property facts: how many concurrent lanes touch it, whether it is a safety or a liveness property, what state it ranges over. From those facts the model \emph{derives} a verification tier --- an invariant crossing multiple concurrent lanes demands exhaustive state-space exploration; a single-lane invariant demands a property test; a temporal obligation demands a model checker run. The tier is not advice: a check walks the model and fails when a mandated checker does not exist on disk. The supervisor does not decide, invariant by invariant, what evidence concurrency deserves. The representation decides, and the environment enforces the decision.

% <!-- point: the-iam-join-derives-grants | The IAM join: service relationships × deployment wiring × runtime identity derive the invoker-grant plan, and a coherence gate asserts the derived grants at deploy time. | terms: alignment-principle -->

\textbf{Service relationships × deployment wiring × runtime identity → an access-grant plan.} The typed inter-service edge model declares which service calls which, over which plane, under which identity. Deployment machinery consumes those edges and derives the invoker grants and token audiences each service needs --- and a coherence gate asserts at deploy time that the granted permissions match the modeled relationships. Least-privilege wiring is not a checklist a human re-audits; it is a projection of the service graph.

% <!-- point: the-provenance-join-derives-the-changelog | The provenance join: every mutation stamps its verb and pass, the customer changelog is derived from the stamps, and a wiring validator catches mutating paths that fail to stamp. | terms: alignment-principle -->

\textbf{Mutation → stamp → derived account.} Every mutating operation in the remediation core writes a provenance stamp --- the verb, the pass, the visibility of the change --- and the customer-facing changelog is \emph{derived} from those stamps rather than authored beside them. A wiring validator identifies mutating paths that fail to stamp their work. The account of what the factory did to a document is therefore evidence, not narrative.

% <!-- pullquote -->

\begin{quote}
Modeling changed what the factory could know. Alignment changed what that knowledge could cause.
\end{quote}

\subsection{The Computation Graph: a Representation That Took Over Execution}\label{heading:the-computation-graph-a-representation-that-took-over-execution}

% <!-- point: the-computations-stabilized-before-composition | The remediation core was decomposed into passes early; their composition stayed implicit until it accumulated global cost. | terms: modeling-principle -->

The computation graph became the deepest of these joins over time. Its history shows why. The remediation core was organized early as passes --- identifiable units of computation --- and that decomposition proved durable. What stayed implicit was their composition: no pass declared, in one common vocabulary, what it produced, what it consumed, what it read only to decide whether to run, and how bounded its mutations were. Across a growing pipeline those locally reasonable freedoms interacted: work whose local complexity looked harmless accumulated quadratic behavior as passes re-derived state, and passes mutated the document through different mechanisms.

% <!-- point: deriving-the-graph-exposed-the-distinctions | Deriving the missing composition exposed distinctions the code had never stated: data flow versus control gating, and the boundedness of each pass's mutations. | terms: modeling-principle -->

The response was to represent the composition well enough to reason about it. Deriving the missing producer--consumer relations from the passes raised a sharper question --- consumer in what sense? Roughly two-thirds of the candidate relationships were not data flow at all but routing: a pass read a verdict only to decide whether it should run. So the passes acquired distinctions they had never expressed uniformly --- produces, consumes, consumes-for-control --- and the graph projects separate data-flow and control-gate relations from them. The same exercise classified every pass site by mutation kind: typed-patch producer, direct editor, or read-only. The counts did not say the direct editors were wrong; they showed where degrees of freedom remained. \autoref{modeling-history} draws the progression.

\begin{figure}[htbp]
\centering
\includegraphics[width=\linewidth]{book/latex/assets/modeling-history-b6b329bc6b.pdf}
\caption{\emph{The modeling history.} DocAble acquired a stable decomposition of remediation computations before it acquired an explicit model of their composition. As locally reasonable choices accumulated consequential global cost, modeling exposed dependency semantics and effect boundedness. Typed declarations made those relationships analyzable and checkable; later redesign made explicit composition authoritative for execution, while separate analytical views continued to grow around stable computation identities. The progression is from a partial model to a richer and increasingly consequential one, not from no model to model.}
\label{modeling-history}
\end{figure}

% <!-- point: the-representation-became-authoritative-for-execution | The pass declarations became the source the graph projects from, a blocking parity check holds them together, and a later redesign made the modeled dependencies authoritative for execution order. | terms: governance-conversion -->

Two further moves made the representation consequential. First, the relationships stopped being a second hand-written truth: the pass declarations carry the typed IO, the graph projects its edges from them, and a blocking parity check fails when declaration and projection disagree. Second, a later redesign made the model authoritative for execution itself: the execution machinery consumes the modeled dependencies to determine when work becomes ready and in what order it may proceed, so composition no longer lives implicitly in the pass bodies. \magevalue{sec:interlude-one-problem-many-models} traces how that architecture emerged from a roughly 250 MB presentation that exceeded the limits of page-chunked execution. In the finished factory, the model determines execution rather than merely describing it. The analytical graph remains as a view over the modeled structure, carrying questions execution need not answer: critical paths, mutation boundedness, latency and cost, test coverage.

% <!-- point: explicit-models-changed-what-was-checkable | Explicit lifecycle and ownership models made previously slippery properties checkable — exhaustive exploration of the ownership state space, temporal specifications checked mechanically — and created the correspondence obligation in return. | terms: modeling-principle -->

Explicit models changed not only what the system recorded but what questions the environment could ask of it. A lifecycle model made illegal transitions explicit. An ownership model reduced a concurrency protocol to the states and interleavings needed to ask whether two workers could hold the same claim; a small purpose-built checker could then exhaustively explore that finite state space. Eventual termination --- a property over executions, not states --- was written as small temporal specifications and checked mechanically; with recovery disabled, the checker produced exactly the stranded execution the property was meant to exclude. Checking the model created a second obligation: why believe a property established over the model holds of the implementation? Conformance tests exercised production operations against the executable model, and when one such test drifted red for days after production changed, the durable response was correspondence machinery --- model-to-implementation checks selected automatically when modeled surfaces change. The factory then added machinery to keep the map true.

\subsection{Exploring Alternative Reasoning Engines}\label{exploring-alternative-reasoning-engines}

% <!-- point: reasoning-substitution-reuses-engineered-operations | Model selection changes the reasoner assigned to a task; the OpenAI pilot changed reasoning orchestration while reusing DocAble's deterministic document operations. -->
% Evidence: docable-deterministic-reasoning-separation-pilot-261029.md; openai-hosted-office-remediation-agent-261029 PHASE3-READOUT.md. Integration and corpus evaluation are distinct stages; shared prompt consumption remains a proposed follow-up.

DocAble's separation of reasoning from deterministic document operations permits substitutions at different levels. For individual tasks, ModelSelector assigns an eligible foundation model using task-specific model tiers and capability requirements. The remediation graph remains unchanged; the model performing a particular judgment can change.

In October 2026, I explored a broader substitution: replacing the graph's reasoning orchestration with a general-purpose OpenAI agent. The agent received a subset of DocAble's engineering models as obligations, leaving it to decide how to accomplish the remediation. Separate skills explained how to inspect documents, make semantic judgments, and invoke existing tools. Those tools supplied the production pipeline's accessibility checks, typed document mutations, and content-preservation validation.

The integration took a few hours and required no changes to DocAble's deterministic C\# code. The agent selected and sequenced operations through existing typed interfaces. This preserved engineering work invested in document representation, mutation, and validation while changing the reasoning strategy. The pilot exercised an alternative to the graph's orchestration; it neither replaced the production pipeline nor established that every computation represented in the graph had an agent-driven counterpart.

An initial trial on five presentations incurred roughly four times the production pipeline's inference expenditure. Quality remained an open engineering problem: two outputs could not be retrieved, retrieved outputs had lower deterministic accessibility scores than production outputs, and an independent semantic evaluation remained outstanding. These results favored the specialized configuration's observed economics without establishing an inherent advantage for graphs. The architectural result was the reuse of engineered document operations under a different reasoning strategy.

\section{The Factory in Operation}\label{heading:the-factory-in-operation}

% <!-- section-terms: governance-conversion, governed-engineering-environment -->

% <!-- point: the-mature-production-flow | The mature production flow in one pass — what the build accumulated, in motion: intent becomes a work unit, realization fans out to isolated worktrees, and admission is serialized through an integration gate that regenerates the evidence against the actual candidate tree. -->

With the representations and joins in place, the mature factory\textquotesingle s production flow can be walked in one pass. What follows is everything the build accumulated, in motion. Engineering intent becomes a named work unit; a brief carries the task, its scope, and the context selected for it; realization fans out to agents working in isolated copies of the repository. Nothing an agent produces reaches the production branch directly. Each candidate change queues for a merge train that rebuilds, re-lints, and re-tests the actual candidate tree --- the independent-evidence discipline a later episode in this subsection explains --- and only work that survives that gate lands. Parallelism lives entirely on the realization side; admission stays serialized through regenerated evidence.

% <!-- point: task-scoped-context-connects-models-to-dispatch | Dynamic context injection supplies task-specific knowledge at dispatch and after planning; its extraction gate requires discovery, while separate controls enforce selected obligations. -->

DocAble selects task-specific engineering knowledge through dynamic context injection (\magevalue{sec:the-agent-stack}). The initial work order supplies an anticipated file set; after planning, the agent repeats the extraction using its refined scope. A pre-commit gate requires evidence that discovery occurred, while separate controls enforce selected obligations. This lets the factory maintain a large modeling substrate without placing all of it into every agent's context.

% <!-- point: repair-is-cheap-escalation-is-the-signal | Most failures die where they are cheap: the agent repairs against the accumulated local gates, a failed admission returns the work for another attempt, and only failures that resist local repair — or implicate the environment itself — reach the engineer. -->

Between realization and admission sits repair, and most failures die there, where they are cheap. An agent whose change fails a local check --- the accumulated controls again, now doing their routine work --- repairs it and tries again; a candidate that fails the merge train\textquotesingle s gate is rejected without landing, and the work returns for another attempt. The engineer sees none of this. What reaches the engineer is the residue: failures that resist local repair, evidence the machinery cannot interpret, and --- the consequential class --- failures that implicate the engineering environment itself. Escalation, in this factory, is not a ticket queue; it is the moment a failure stops being an implementation problem and becomes an engineering one.

\subsection{What Escalated Failures Became}\label{heading:what-escalated-failures-became}

% <!-- point: governance-conversion-is-an-empirical-claim-about-durable-structure | Where did the machinery come from? Not grand design: governance conversion claims some episodes leave durable structure later work inherits, and this subsection reconstructs the episodes where that sequence is visible — the accumulation the build's control curve counted in aggregate, seen one episode at a time. | terms: governance-conversion -->

Where did all this factory machinery come from? Not from a grand upfront design. The build section traced the accumulation of models and checks; what follows reconstructs where selected pieces of it came from, one episode at a time. Governance conversion makes an empirical claim: some engineering episodes leave behind durable structure that later work inherits instead of paying for the same judgment again. What follows reconstructs selected DocAble incidents where that sequence --- pressure, response, durable consequence --- is actually visible, organized by what each failure \emph{converted into} rather than by when it happened. It is not a claim that every structure arose from failure; the subsection ends with the ones that did not.

% <!-- point: implicit-lifecycle-became-explicit-model | An implicit lifecycle and a retry reflex became an explicit lifecycle model, atomic claiming, and recovery as a decision over evidence. | terms: governance-conversion -->

\textbf{An implicit lifecycle became an explicit behavioral model.} The dominant production failure was a slow generation call, and the lifecycle around it was implicit --- status written from scattered places, one terminal state masking many distinct failures, and requeue as the reflex recovery. Under a slow dependency that reflex became a storm: requeue, still slow, same timeout, requeue again --- the recovery amplifying load exactly when the system was already hurting, every attempt billing its full cost with nothing marking it a retry. What became durable was the explicit lifecycle model from the previous subsection: named states and legal transitions, atomic claiming that separates permission to act from recovery leases, and no automatic requeue on ordinary failure. A later stale-owner sweep exposed the next missing distinction --- slow is not dead --- and progress evidence supplied it. Recovery stopped being a reflex and became a decision over explicit state and evidence.

% <!-- point: destructive-library-behavior-became-a-typed-seam | A convenient library call that silently destroyed accessibility structure became a canonical-library policy, a typed mutation seam, and a build-time ban — and the gate itself later had to be hardened. | terms: governance-conversion -->

\textbf{Destructive library behavior became a typed seam and a ban.} An early path used a convenient scripting-language library to slice a document. The operation worked visually and silently destroyed the internal tag structure carrying the accessibility the product exists to add --- and nothing caught it. The response came in layers: a policy named the canonical library; a typed format model and API concentrated all mutation of that format through one seam; a build-time check refused raw access from elsewhere. Then the check itself failed --- a stale repository root caused real findings to be reported as green --- and the second-order incident forced the gate to be hardened too. The model stated the boundary; the check enforced it; the enforcement machinery itself still had to be maintained.

% <!-- point: cross-service-assumptions-became-explicit-models | A file path sent across a service boundary became a typed client that makes the mistake unrepresentable, plus an explicit service graph — and a review found the bypass path the boundary did not govern. | terms: governance-conversion -->

\textbf{A cross-service assumption became a typed client and an explicit service graph.} Another failure sent a local file path across a service boundary, where it became meaningless in the callee\textquotesingle s process. The replacement client accepted file content but not a local path --- the mistake became unrepresentable at the type level --- and an explicit service graph represented which services could call which, the same graph the access-grant join now derives from. A later human review found a second problem: an optional direct-client path could bypass shared rate limiting, key rotation, failover, and cost tracking entirely. A protection boundary only governs the paths that cross it.

% <!-- point: an-instruction-satisfied-to-the-letter | "Use types," satisfied to the letter and voided in spirit, became a commodity-lint floor driven to zero plus induced typed abstractions region by region. | terms: stringly-typed, governance-conversion -->

\textbf{An underspecified instruction became a lint floor and named abstractions.} By the time the system reached roughly 300,000 lines, regions built MVP-fast lacked coherent abstractions, so I told the agent to use typed Python. It did, in a sense: it used the type \texttt{string} for everything --- stringly-typed, the word "type" satisfied to the letter and voided in spirit. "Use types" specified a surface property, not the domain abstractions the types were supposed to represent, and the agent took the cheapest satisfying path. What made the period alarming was the speed with which a weakly specified direction could turn into a large amount of coherent wrong structure. The durable response combined a floor with structure: every commodity lint the language ecosystems ship was turned on and driven to zero, each made blocking once the repository conformed; then, region by region, dense primitive code was routed through induced typed domain abstractions, with the checks promoted as each governed surface reached conformance. A month later the 300,000 lines were roughly 200,000, and future work inherited named abstractions, boundaries, and regression checks rather than a rewritten snapshot.

% <!-- point: optimizing-within-an-architecture-that-was-itself-the-constraint | Weeks of in-frame autoscaler optimization against the same ceiling became a decision rule: price the architectural change instead of commissioning another local optimization. | terms: velocity-exposes-the-danger, governance-conversion -->

% <!-- index-def: velocity-exposes-the-danger -->

\textbf{A recurring architectural ceiling became a decision rule.} The service cost tens of dollars a day while idle on a cluster. The first response stayed inside the existing architecture --- tune the autoscaler, reduce idle pods --- and for weeks the agents made it better and better, with every improvement colliding with the same boundary: cold pods took too long to return. The repeated failure was information about the \emph{frame} rather than the implementation inside it. The durable lesson was a decision rule --- when repeated in-frame optimization collides with the same architectural ceiling, price the architectural change rather than commissioning another round of local optimization --- and the eventual serverless migration encoded the decision in the system itself. A judgment becomes capital only when later work can inherit it; a lesson held in someone\textquotesingle s head remains expensive to reuse.

% <!-- point: an-unmade-decision-is-where-the-agent-guesses-worst | A conflict handler with an unmade design decision silently deleted landed work for four days; the choice was made once, encoded, and the destructive default barred. | terms: optionality-is-poison, governance-conversion -->

% <!-- index-def: optionality-is-poison -->

\textbf{An unmade decision became a decision made once.} A merge automation kept stalling when two agents edited the same file, and the proposed one-flag fix --- on conflict, take one side and drop the other --- was signed off without asking what "drop the other side" meant. For four days it silently deleted real work from the main branch, the same commit heading landing again and again, each time a little smaller. The conflict handler contained a consequential design choice that had never been made explicitly: what information may be discarded when two changes overlap? At execution time the agent still had to choose, so it selected the locally expedient behavior. The resolver was changed to reconcile the conflicting files, and the destructive default was prohibited. The portable lesson is narrow: do not leave consequential semantics unresolved at a boundary where an autonomous actor must act.

% <!-- point: two-green-reports-that-were-not-true | Two green reports — a lint gate swallowing its own error, pin tests landed onto a tree missing the checked field — became independent evidence regenerated at the admission boundary. | terms: done-is-a-claim, governance-conversion -->

% <!-- index-def: done-is-a-claim -->

\textbf{"Done" became a claim requiring independent evidence.} Two incidents, close together, made the rule non-negotiable. A lint gate hit a stray error, caught it in a catch-all that returned success, and let a run of commits land completely unlinted while the log said fine. A chain of pin tests passed green --- because a botched merge had landed them onto a tree missing the very field they were supposed to check. Both reported done; neither was. A report of "done" describes the producer\textquotesingle s claim, not the state of the repository, and the gap between the two is where a confident green hides. The durable response regenerated the evidence at the admission boundary: required tests and lints re-run against the actual candidate tree, gate failures made fail-loud --- a control cannot translate its own uncertainty into success --- and counts re-derived rather than accepted from the task report. The capital is not a standing habit of distrust; it is machinery that regenerates the evidence cheaply. That machinery is the merge train whose weekly funnel opened this subsection --- the same gate, measured in operation.

% <!-- point: a-detector-that-measured-the-right-quantity-on-the-wrong-world | A content-free image detector, provably safe on the world it was measured on, would silence a Rothko; the cheap proxy lost the last word and a regression preserved the failure class. | terms: tests-for-agent-failure-modes, governance-conversion -->

% <!-- index-def: tests-for-agent-failure-modes -->

\textbf{An inadequate input representation became a routed judgment.} An agent wrote a clean, typed, careful detector: if an image\textquotesingle s pixels barely vary, skip the vision model and mark it decorative. Its own comment promised a false positive was "not possible" --- real photos and charts vary twenty times more than the cutoff. The margin was true, and it was measured on the wrong world. Hand it a Rothko, or Malevich\textquotesingle s black square, and the detector marks the painting decorative, and a blind student reading an art book gets silence where the subject should be. No threshold saves it, because a blank spacer and a low-variance artwork can be identical in the one quantity the code measures --- "content-free" is a judgment about an image in context, not a fact about pixels. The cheap proxy lost the last word: ambiguous cases route to a richer semantic judgment, and the regression evidence covers the failure class --- low-variance images that nevertheless carry content --- rather than the first example that exposed it.

% <!-- point: bypasses-produced-empty-changelog-rows | Eleven mutation sites bypassing the stamp helper shipped blank changelogs; representation, migration, then gating turned provenance from inference into a lookup. | terms: governance-conversion -->

\textbf{Forensic reconstruction became derived provenance.} A corrupted document once took four cross-file hops to trace back to the pass that wrote it, and verification then found eleven production sites writing through a low-level path that skipped the provenance stamp --- shipping empty changelog rows for a whole class of jobs while the code had changed plenty. The response is the provenance join described earlier. Its landing sequence matters: the wiring validator \emph{reported} first, because the inherited tree contained many violations; a fix wave drained the gap; only then was the obligation enforced at the admission boundary. Representation made the obligation observable; migration made it feasible; gating enforced it.

% <!-- point: a-timeout-spoke-only-after-the-budget-was-spent | A silent eight-minute stall became a provisional per-chunk worst-case model that is deliberately NOT a gate — alignment need not culminate in blocking enforcement. | terms: governance-conversion -->

\textbf{A silent stall became measurement --- and deliberately stopped there.} A document\textquotesingle s repair stalled in silence; the subprocess was killed by the job timeout after roughly eight minutes, the customer\textquotesingle s progress bar frozen at thirty-eight percent, the budget already spent. The incident produced a provisional per-chunk worst-case model against the job budget --- the seed kept as a timestamped datapoint in a file, not frozen as a constant in code --- and instrumentation that can report a would-be breach. \autoref{measurement-no-gate} draws where the chain deliberately stops: no production gate rejects a job on this model today, because the bound is not calibrated strongly enough to justify enforcement, and a false refusal could cost more than continued observation. \textbf{Alignment need not culminate in a blocking gate.} The counterpart shows measurement earning action when its evidence is adequate: a deterministically measured per-request cold-start floor of 4,057 ms, joined to a deployment-topology model, drove a re-architecture that cut the floor to 109 ms. Same progression from representation to evidence; different verdicts, each proportionate to the evidence.

\begin{figure}[htbp]
\centering
\includegraphics[width=\linewidth]{book/latex/assets/measurement-no-gate-300b62ff05.pdf}
\caption{\emph{Modeled and Observed, Not Yet Binding.} A timeout incident produced a measurement, a provisional per-chunk model, and reportable evidence. Production admission still does not depend on the model: its bound is not calibrated strongly enough to justify blocking. The terminal enforcement node is left open on purpose --- representation and observation can mature before enforcement is earned. The open node is part of the evidence, not an omission.}
\label{measurement-no-gate}
\end{figure}

% <!-- point: some-conversions-were-almost-mechanical | Some conversions were mechanical: a repeated review finding — a missing reserved prefix on inserted artifacts — became a blocking lint. | terms: governance-conversion -->

\textbf{And some conversions were almost mechanical.} Invisible remediation artifacts were supposed to carry a reserved name prefix so a validator could always tell an inserted element from an authored one, and review kept catching insertions that skipped it. A prefix rule is entirely decidable, so the audit question became a blocking lint: the judgment moved from repeated inspection into the environment, where it fires the same way on every commit.

% <!-- point: the-case-contains-both-routes-to-durable-structure | Not every structure was forced by failure: the reactive backend, the closed edit vocabulary, and the unified rule engine were designed from known obligations and simply held. | terms: governance-conversion -->

\textbf{Not every structure was purchased by an incident.} The backend became reactive and stateless early, which later made the serverless migration cheap; the document editor routed human gestures and automated remediation through one closed edit vocabulary, which later supported a second producer without a second mutation path; the unified detect-and-repair engine consolidated two scanning paths for a clean architectural reason, and a later audit found the boundary holding. Report those as confirmation, not conversion. Some assets arose from incidents and were refined through repeated use; others were designed from known obligations and proved effective as the system evolved. The factory\textquotesingle s machinery accumulated along both paths, without a comprehensive plan established in advance.

\section{The Supervisor\textquotesingle s Job}\label{heading:the-supervisors-job}

% <!-- section-terms: governance-conversion, governed-engineering-environment -->

% <!-- point: the-preceding-subsections-recorded-substitutions | The preceding subsections each already recorded a piece of the answer; the staircase names the same transformation chronologically, from the supervisor's chair. It is a chronology of DocAble, not a maturity model. | terms: governance-conversion -->

What was the human doing while the factory accumulated? The engineer\textquotesingle s responsibilities changed as more work became delegated. The preceding subsections have each already recorded a piece of it; the staircase below names the same transformation chronologically, from the supervisor\textquotesingle s chair: each step moved some combination of system knowledge, realization work, evaluation, or admission out of direct human handling and into the engineered environment. Five stages name that progression in my own work: co-coder, QA, review lead, tech lead, and architect. What matters is what the engineered environment could carry at each stage. It is a chronology of DocAble, not a maturity model for other factories.

% <!-- point: rungs-compressed | Each rung delegated what the environment had just become able to carry: reading, then evaluation, then decidable findings, then locality of reasoning, then system-level knowledge itself. | terms: governance-conversion -->

\textbf{Co-coder:} direct human inspection carried everything --- I read the agent\textquotesingle s changes, judged them, merged them. It worked while volume stayed inside one person\textquotesingle s attention, and failed as an operating model once generation outran review. \textbf{QA:} evaluation was delegated before enforcement --- a second agent audited another agent\textquotesingle s work, episodically and probabilistically; its value was diagnostic, and its repeated decidable findings revealed questions that no longer needed a language model at all. \textbf{Review lead:} reviewers specialized by concern and region, and every finding that turned out to be mechanically decidable dropped into deterministic machinery, the probabilistic reviewers keeping the questions checks could not settle. \textbf{Tech lead:} representation made reasoning local --- named zones externalized ownership and boundaries so an agent could reason inside one region without first recovering the whole system, and selected obligations expressed over the zone model were enforced. \textbf{Architect:} by the final rung I rarely read implementation directly. Architectural, behavioral, ownership, and other models supplied the system views for planning and review; agents used the same representations to localize changes. Declaration did not make the models true --- model drift became its own failure class, which is why correspondence checking earned its place in the previous subsections. And at this rung the economics of restructuring inverted: a cross-format migration that touched \magevalue{refactor_files} files and rewrote \magevalue{refactor_lines} lines took about \magevalue{refactor_hours} hours of agent time over a dinner break. The hard part was never touching \magevalue{refactor_files} files; it was knowing those files should be touched. \autoref{mage-staircase} draws the climb.

\begin{figure}[htbp]
\centering
\includegraphics[width=\linewidth]{book/latex/assets/mage-staircase-ad5d286bc2.pdf}
\caption{\emph{The Delegation Staircase.} As explicit representation and environmental enforcement accumulated, larger units of engineering work could be delegated. Representation moved system knowledge out of one person\textquotesingle s head; enforcement moved repeatable admission decisions out of direct review. Human work moved from reading every diff toward system-level strategy and the residual semantic judgment the environment could not yet decide. A failure in either representation or enforcement could move work back down the staircase.}
\label{mage-staircase}
\end{figure}

% <!-- point: the-staircase-has-two-interacting-supports | The staircase's two supports: Alignment moved enforcement of repeatable obligations out of direct review; Modeling moved system knowledge out of the engineer's head; weakness in either sent the work back down. | terms: alignment-principle, modeling-principle -->

The staircase has two interacting supports, and delegation climbed only as far as they could carry it. \textbf{Alignment} moved enforcement of repeatable obligations out of direct human review: deterministic checks, permissions, boundaries, and gates enforced obligations I no longer inspected line by line. \textbf{Modeling} moved system knowledge out of my head: named zones and later explicit models let agents reason over larger engineering questions without reconstructing the whole repository each time. When representation was too weak, I fell back into reconstructing the map; when enforcement was too weak or unreliable, I fell back into reading diffs and verifying outcomes by hand. A miss sent the staircase backward.

% <!-- point: human-attention-moved-to-a-functional-list | Functionally, the top of the climb rewrites the supervisor's job description: human attention moved away from supervising individual acts of realization and toward the work only judgment can do. | terms: governed-engineering-environment -->

Functionally, the top of the climb rewrites the supervisor\textquotesingle s job description. As the MAGE layer strengthened, human attention moved away from supervising individual acts of realization and toward:

\begin{itemize}
\tightlist
\item
  deciding which external pressures become engineering work;
\item
  choosing which engineering questions are consequential enough to represent;
\item
  constructing and revising the representations themselves;
\item
  deciding which obligations deserve enforcement, and at what strength;
\item
  interpreting anomalous evidence the machinery surfaces;
\item
  resolving genuinely ambiguous tradeoffs;
\item
  connecting representations to one another --- the joins of the earlier subsection;
\item
  diagnosing failures in the engineered environment itself;
\item
  and deciding what experience should become durable structure.
\end{itemize}

% <!-- point: prompt-properties-preserved-useful-freedom | Prompt failures exposed evidence-grounding and prefix-ordering obligations; representing and checking these properties preserved freedom over prompt wording. -->

Prompt design exposed two consequential properties I had initially left to the agents. Agents writing prompts for DocAble's remediation models imposed numerical demands, such as a minimum number of facts in an image description. When the available evidence could not support that many facts, the demand encouraged invention. They also placed variable document context before fixed instructions, losing opportunities to reuse the model's key-value cache. On one 70-slide presentation initially costing approximately \$1.50, moving fixed content to the prefix reduced inference cost approximately threefold.\citep{davis2026docable}

I responded by modeling prompt properties by kind and adding checks. Requiring descriptions to be grounded in source evidence addressed the semantic problem more directly than replacing one numerical threshold with another. Requiring stable content to precede variable context addressed the economic problem without prescribing the prompt's wording. The remaining wording choices could stay with the agents. Neither response required me to write every prompt; both required deciding which properties of a prompt had consequences worth governing.

% <!-- point: mage-afforded-supervision-by-changing-its-objects | MAGE afforded supervision by changing the objects through which supervision occurred — the section's observed substitutions named as one claim. It did not automate supervision, and human judgment remains central. | terms: governed-engineering-environment -->

Implementation gave way to engineering intent as the unit of supervised work; reading source gave way to explicit representations of consequential properties; the producer\textquotesingle s claim gave way to independently regenerated evidence; routine inspection gave way to anomalies and escalations; per-diff review gave way to serialized admission decisions; isolated descriptions gave way to joins that carry consequences; and recurring local judgment was converted, episode by episode, into durable structure. \textbf{MAGE afforded supervision by changing the objects through which supervision occurred.} Not "MAGE automated supervision" --- human judgment remains central, and every item on the list is judgment. But supervisory effort no longer scales linearly with implementation volume for any property that can be represented and mechanized, and that is the difference between the co-coder reading every diff and the architect reading almost none.

\section{Economics}\label{heading:economics}

% <!-- section-terms: remediation-cost -->
% <!-- point: the-accounting-prices-the-case | Construction, operation, and retained human review are different costs; the SEIP development estimate is author-supplied. -->

A factory's startup cost, its marginal change cost, and its carrying cost are different numbers --- a separation \magevalue{sec:toward-a-theory-of-mage} develops formally. The SEIP accounting estimates approximately \$110,000 to develop DocAble over six months: \$90,000 in engineering labor and \$20,000 in cloud expenditure. These are author-supplied estimates, not an independently audited ledger.\citep{davis2026docable} Much of that expenditure purchased reusable capability. Dividing it by the documents processed during development would confuse the cost of building the factory with the cost of using it.

\subsection{Human Remediation as a Reference}\label{heading:human-remediation-as-a-reference}

% <!-- point: manual-remediation-provides-a-reference-not-an-oracle | The undergraduate corpus provides actual artifacts and billed labor; it does not establish expert ground truth or matched per-document labor savings. -->

The undergraduate effort introduced in \magevalue{sec:mage-by-example} supplies a concrete reference. Across approximately six months, \magevalue{seip_uta_workers} trained undergraduates billed \magevalue{seip_uta_hours} hours, and the recovered corpus contains \magevalue{seip_uta_pairs} identifiable document pairs from \magevalue{seip_uta_courses} courses. The resulting \magevalue{seip_uta_hours_per_pair} hours per pair is a corpus-level ratio, not a stopwatch measurement for each document. The students' outputs are evidence of an institutional remediation workflow; they are not an expert oracle for accessibility.

I developed DocAble using independently obtained documents and the accessibility standards, and withheld examination of the students' remediated files until near the end of development. This reduces one route to contamination without establishing complete blinding. One spreadsheet in the validation partition had prior development exposure.

The recovered student outputs provided a reference for evaluating DocAble's remediation. Eligibility screening and a format-stratified split produced 63 validation documents; DocAble completed 62, failing on one PDF. In a blinded, transcript-based GPT-5.2 rubric assessment of those outputs, 20 scored better than the student output, 41 fell within a one-point band, and one scored worse. The rubric does not assess all visual properties or screen-reader usability, and Microsoft changed its checker after the students' remediation work. These findings support further evaluation, not a claim of complete accessibility.

Median runtime across all 63 validation attempts was 1.4 minutes; median inference cost for the 62 completed outputs was \$0.066. Those quantities omit retained human review and do not include a complete cloud-cost accounting. Comparing billed human hours directly with unattended machine runtime would overstate what we know about labor savings. The practical question is how much competent review and correction remain after either workflow.

\subsection{General-Purpose Agents and a Specialized Service}\label{heading:agents-and-specialized-service}

% <!-- point: agent-comparison-tests-economics-without-a-capability-boundary | A small comparison finds mixed quality and lower observed runtime and inference cost for DocAble; rubric acquisition and delivery remain engineering concerns. -->

Could a general-purpose agent perform the remediation directly? The SEIP comparison uses community document tools, unlike the DocAble tools reused in \hyperref[exploring-alternative-reasoning-engines]{Exploring Alternative Reasoning Engines}. It tests Opus and Fable on ten validation documents: five PDFs, two Word documents, two presentations, and one spreadsheet. Each agent used one configuration. \autoref{agent-remediation-comparison} summarizes completion, runtime, and inference cost for this ten-document cohort. Its DocAble cost figure concerns this smaller cohort, rather than the 62 completed outputs underlying the \$0.066 median above.

{\footnotesize
\begin{longtable}[]{@{}lrrr@{}}
\caption[Agent and service remediation comparison]{Comparison of three remediation configurations on ten documents. Medians cover completed attempts; Fable timed out on one PDF. Agent inference costs are cache-aware API-equivalent estimates excluding cloud execution and evaluation.}\label{agent-remediation-comparison}\tabularnewline
\toprule
\textbf{Method} & \textbf{Completed} & \textbf{Minutes} & \textbf{Inference cost} \\
\midrule
\endfirsthead
\toprule
\textbf{Method} & \textbf{Completed} & \textbf{Minutes} & \textbf{Inference cost} \\
\midrule
\endhead
\bottomrule
\endlastfoot
Opus & 10/10 & 8.22 & \$2.466 \\
Fable & 9/10 & 7.77 & \$3.977 \\
DocAble & 10/10 & 1.31 & Less than \$0.50 \\
\end{longtable}
}

Quality varied by format. On the four PDFs completed by every condition, median reductions in veraPDF findings were 89.7\% for Opus, 91.8\% for Fable, 100\% for DocAble, and 75.3\% for the student outputs. Both agents outperformed DocAble on Office checks for the two Word documents; the student outputs achieved the largest reductions for the two presentations. The spreadsheet had no original findings. Checker reductions measure selected properties, so neither the rankings nor completion establish overall accessibility.

A capable general-purpose agent supplied with a sufficiently complete rubric and suitable tools could match or exceed DocAble. The engineering problem includes acquiring that rubric: experience with difficult documents reveals omissions that a plausible initial specification can miss. Some of this knowledge may become available through model training; some must be supplied as task-specific context or encoded in tools and checks. A specialized service packages accumulated knowledge and the workflow for users, and can reduce the cost of applying it repeatedly. Its advantage depends on the quality of that knowledge, operating economics, and usability. The comparison measures particular configurations at a particular time; it establishes neither a lasting capability boundary nor a causal benefit from MAGE.

\subsection{What the Economics Establish}\label{heading:what-the-economics-establish}

% <!-- point: measured-production-cost-leaves-review-and-carrying-cost-open | The evidence separates startup investment, inference expenditure, and remaining human work; it does not establish institutional savings or a market size. -->

The three alternatives incur costs in different places. Manual remediation purchases labor for each document. A general-purpose agent needs a useful rubric, appropriate tools, and supervision. A specialized service incurs engineering and maintenance costs to make the recurring workflow easier and cheaper to run.

The case makes a narrower economic claim than a university-wide savings estimate. Producing remediation outputs can require little inference expenditure once the supporting system exists. Establishing that those outputs serve their readers, integrating the service into institutional work, and maintaining it remain costs to measure. Software can distribute the resulting capability across courses and institutions, but replication does not eliminate review, integration, or upkeep. Those are part of evaluating the factory's return, rather than exceptions to it.

\section{Liability}\label{heading:liability}

% <!-- section-terms: governance-conversion -->

% <!-- point: the-map-pointed-at-a-ghost | A model kept naming a moved symbol and stayed authoritative-looking for days — accumulated structure fails in its own ways. | terms: governance-conversion -->

The factory also introduced substantial maintenance costs and new failure modes. One model named a code symbol that was later moved and renamed. Nothing invalidated the reference; the model continued to look authoritative, agents continued to reason from it, and the failure surfaced days later during deployment. A later audit showed the stale pointer was a class, not a one-off --- re-running closed work against current code found genuine model-to-code drifts a green definition of done had missed. That incident is why correspondence checking exists, and it generalizes.

% <!-- point: accumulated-structure-has-carrying-costs | The liability list: models drift, controls conflict, evidence goes stale, gates encode obsolete assumptions, the support apparatus needs operating, and some of what accumulated stops being worth its upkeep. | terms: governance-conversion -->

The models can drift, and the correspondence machinery that catches that drift must itself be maintained. Controls can conflict, and resolving a collision between two well-intentioned gates is supervisory work no model retires. Evidence goes stale: a baseline recorded under last quarter\textquotesingle s workload quietly stops meaning what its readers assume. Gates encode the assumptions of the incident that created them, and an obsolete assumption enforced mechanically is worse than one merely remembered. The support apparatus, about 3.8 times the production code at the final snapshot, must itself be operated, debugged, and paid for in attention. The structures that carry decisions forward must themselves be maintained, and some will stop being worth their upkeep: not all of what this factory accumulated will survive its next architectural turn.

% <!-- point: the-claim-is-bounded-one-engineer-one-ecosystem | The evidence has clear limits: one longitudinal case, one primary engineer, one model ecosystem, no controlled comparison — the case establishes mechanism and sequence, not prevalence. | terms: field-report -->

The evidence has clear limits. This is one longitudinal production case led by one primary engineer directing an agent fleet, using one contemporary model ecosystem, without a controlled comparison against another engineering process. It cannot establish that another organization will encounter the same failures, build the same mechanisms, or achieve the same economics --- and checker scores do not establish complete end-user accessibility. What the case can establish is mechanism and sequence: these pressures occurred, these responses followed, and selected structures persisted long enough to be exercised under continuing change.

% <!-- point: turn-outward-to-other-factories | The transition: DocAble is one deliberately engineered response to abundant realization, so a factory has to be evaluated as a production system rather than by agent capability alone; the next question is whether independently developed factories exhibit the same pressures. -->

DocAble shows one deliberately engineered response to abundant realization. The factory therefore has to be evaluated as a production system, including the structures that direct and supervise realization, rather than by agent capability alone. The next question is whether independently developed agentic software factories exhibit the same pressures --- and that question does not depend on anyone else having heard of MAGE.
